% Source tables: T3, T4
\begin{table*}[t]
\centering
\small
\setlength{\tabcolsep}{4pt}
\caption{Internal design comparisons on native VisDrone validation results.}
\label{tab:selected-3}
\textbf{(a) Localization design}\par\smallskip
\begin{tabular}{lrrrr}
\toprule
Variant & Geometry & LRS & AP50 & AP \\
\midrule
Original FDR & reference & reference & 48.473 & 29.078 \\
Clean-FDR & -- & -- & 49.320 & 29.696 \\
Clean-FDR + geometry & + & -- & 49.130 & 29.361 \\
LRS-GFDR & + & + & 49.382 & 29.620 \\
\bottomrule
\end{tabular}
\par\smallskip
\begin{minipage}{0.97\textwidth}\footnotesize These are registered design variants, not a fully source-matched toggle experiment. Geometry reduces validation AP relative to Clean-FDR in these runs; LRS recovers part of that difference. The table does not establish a standalone accuracy gain from geometry or prove a second-dataset stability claim. Reference denotes switches not mapped to the current implementation.\end{minipage}
\par\medskip
\textbf{(b) BPDD recipe and FIA compatibility}\par\smallskip
\begin{tabular}{lrrr}
\toprule
BPDD recipe & FIA & AP50 & Best AP \\
\midrule
Capacity v2 & -- & 49.620 & 29.848 \\
Capacity v2 & + & 49.585 & 29.736 \\
Capacity v3-Safe & -- & 48.783 & 29.277 \\
Capacity v3-Safe & + & 49.890 & 30.109 \\
\bottomrule
\end{tabular}
\par\smallskip
\begin{minipage}{0.97\textwidth}\footnotesize The v2-to-v3-Safe change is a recipe bundle, not an isolated toggle for each distillation or scheduling choice. The v3-Safe BPDD-only run is worse than v2, while its joint run is better. These results motivate compatibility-oriented design rather than universally improved standalone BPDD.\end{minipage}
\par\medskip
\end{table*}
